% ================= 3. MARCO REGULATORIO =================
\section{El marco regulatorio: por qué la soberanía es un requisito, no una opción}\label{sec:marco}

\subsection{La nueva LFPDPPP (2025) y el tratamiento automatizado}

El 20 de marzo de 2025 se publicó en el DOF la nueva Ley Federal de Protección de Datos Personales en Posesión de los Particulares, que abrogó la ley de 2010 y entró en vigor al día siguiente.\cite{lfpdppp} El cambio institucional más visible es la extinción del INAI ---por la reforma constitucional de simplificación orgánica de diciembre de 2024--- y la transferencia de sus atribuciones sobre datos del sector privado a la \textbf{Secretaría Anticorrupción y Buen Gobierno (SABG)}.\cite{ccn2025,lfpdppp} Las multas alcanzan \textbf{320,000 UMA} ($\sim\$37.5$ millones de pesos) y pueden duplicarse para datos personales sensibles.\cite{kyc2026}

Tres innovaciones de la ley son directamente relevantes para un sistema de IA jurídica. Primera, la ampliación de la definición de responsable: cualquier persona que realice tratamiento de datos tiene obligaciones directas, lo que captura a encargados y proveedores tecnológicos.\cite{cosio2026} Segunda, la eliminación de los ``fines análogos'': el responsable ya no puede tratar datos para finalidades distintas de las declaradas en el aviso de privacidad, aunque sean compatibles.\cite{ccn2025} Tercera ---y central---, las disposiciones sobre \textbf{tratamiento automatizado}: quien opere decisiones o tratamientos automatizados debe informar su existencia, explicar la lógica general del tratamiento y sus consecuencias posibles, y garantizar el derecho del titular a oponerse cuando produzcan efectos jurídicos adversos.\cite{cosio2026} Un asistente jurídico conversacional es, por definición, un tratamiento automatizado: su arquitectura de consentimiento, abstención y explicabilidad debe diseñarse alrededor de estos artículos, no como añadido cosmético. Al cierre de esta edición el reglamento de la ley seguía pendiente, lo que añade incertidumbre operativa pero no suspende las obligaciones vigentes.\cite{cosio2026}

\subsection{Secreto profesional con sanción penal}

El deber de sigilo del abogado no es solo deontológico. El Código Penal Federal sanciona la revelación de secretos o comunicaciones reservadas conocidas con motivo del empleo, cargo o puesto (art.~210, treinta a doscientas jornadas de trabajo comunitario) y la agrava cuando el revelador presta servicios profesionales o técnicos (art.~211, uno a cinco años, multa y suspensión de profesión de dos meses a un año).\cite{cpf} Enviar el expediente de un cliente a una API extranjera sin garantías contractuales y técnicas de no retención es una exposición penal y de responsabilidad civil que ningún comité de riesgos de despacho puede aceptar de forma permanente. La arquitectura de AI Justicia invierte la carga de la prueba: en el tier de despacho, \textbf{ningún dato sale del perímetro del cliente} ---modelo soberano on-premise, adaptadores LoRA privados, cero llamadas externas--- y el invariante de doble adaptador (Sección~\ref{sec:entrenamiento}) es la expresión arquitectónica de estos artículos.

\subsection{Base normativa del corpus: los textos del Estado no son objeto de derecho de autor}

La viabilidad legal del corpus descansa sobre una regla clara: el art.~14, fracc.~VII de la Ley Federal del Derecho de Autor establece que \textbf{no son objeto de protección} ``los textos legislativos, reglamentarios, administrativos o judiciales, así como sus traducciones oficiales'', añadiendo que su publicación deberá apegarse al texto oficial y no confiere derecho exclusivo de edición.\cite{lfda} Legislación federal y estatal, gacetas oficiales y sentencias judiciales son, por tanto, reutilizables para entrenamiento sin licencia autoral. La doctrina universitaria sigue reglas distintas: las fuentes incorporadas provienen de repositorios institucionales de acceso abierto (IIJ-UNAM vía OAI-PMH, repositorio de tesis UNAM), cuyos términos permiten consulta y reutilización académica; la redistribución del corpus sigue los términos de acceso público de cada fuente, todas ellas gubernamentales o universitarias de acceso abierto.

\subsection{Infraestructura: residencia de datos resuelta, soberanía de inferencia pendiente}

México vive una transformación de infraestructura cloud sin precedente: AWS invirtió más de USD 5,000M en su región de Querétaro (lanzada en 2025)\cite{aws2024} y Google Cloud inauguró la suya en diciembre de 2024, con latencias menores a 5 ms para el centro del país y una contribución estimada de USD 11,000M al PIB.\cite{gcp2024} Esto resuelve la \emph{residencia} del dato ---relevante para los tiers ciudadano y profesional--- pero no sustituye la soberanía de inferencia: el servicio gestionado de un hiperescalador extranjero sigue sujeto a políticas de retención, revisión y jurisdicción ajenas. La taxonomía de soberanía de IA aplicada en la literatura distingue corpus, entrenamiento, capacidad lingüística e inferencia;\cite{alia2026} AI Justicia cubre las cuatro: corpus nacional (Sección~\ref{sec:corpus}), entrenamiento propio sobre pesos abiertos (Sección~\ref{sec:entrenamiento}), capacidad jurídica en español mexicano, e inferencia dentro del perímetro del cliente en el tier de despacho.
